\begin{resumo}[Abstract]
\begin{otherlanguage*}{english}
Managing sales across multiple marketplaces requires repeated product registration and the monitoring of listings, orders, and inventory in separate environments, increasing operational work and the risk of inconsistencies. This work presents the proposal and initial implementation of \sistema{}, a web application intended to centralize these activities, based on the needs of Tutymotos, a business in the motorcycle parts sector. The proposed solution includes product and listing management, order tracking, inventory control, and the analysis of business performance indicators, subject to data availability on the integrated platforms. It also includes artificial intelligence to assist with content preparation, subject to the retailer's review and approval. The procedures adopted include discussions with the administrators, observation of the store's daily operations, requirements elicitation, domain and data modeling, and the organization of the application into layers. Angular was adopted for the web interface, Go for the API, and PostgreSQL for persistence, with data access code generated by sqlc. The outcomes of this stage comprise the initial requirements, system models, and the implementation of an email and password verification flow, which does not yet include session management. All six automated interface tests, using simulated HTTP responses, passed, and static analysis of the API reported no issues. These results do not establish that the system works with the database in an integrated environment or that it provides operational gains for the store. The current contribution consists of defining the problem and building a technical foundation for the system. Assessing its effectiveness requires completing the business modules and marketplace integrations, as well as conducting tests and measurements with the retailers.

\vspace{\onelineskip}
\noindent
\textbf{Keywords}: Marketplaces. Systems integration. E-commerce. Listing management.
\end{otherlanguage*}
\end{resumo}
